\section{天津多联机及风机盘管系统设计}

本章以天津（寒冷地区）为目标城市。天津属于建筑热工设计气候分区中的寒冷B区，冬季较长且干冷，夏季炎热。该城市是本文中代表华北沿海寒冷气候的典型样本。

\subsection{室外设计参数}

天津夏季室外气象设计参数（CSWD气象年数据）为：夏季空调室外干球计算温度\SI{33.9}{\degreeCelsius}，夏季空调室外湿球计算温度\SI{26.9}{\degreeCelsius}，夏季空调室外计算日平均温度\SI{29.4}{\degreeCelsius}，夏季室外平均风速\SI{2.3}{\meter\per\second}。冬季空调室外干球计算温度为\SI{-9.6}{\degreeCelsius}。天津夏季干球温度在五城市中处于中间偏高水平，冬季温度则介于严寒区沈阳与南方城市之间，呈现典型的“冷热并重”过渡型气候特征。

\subsection{建筑概况}

同第一章详述。

\subsection{理想负荷计算结果}

\begin{table}[tbp]
\centering
\caption{天津负荷计算与系统方案比较结果}
\label{tab:tianjin-ideal}
\scriptsize
\begin{minipage}[t]{0.42\textwidth}
\centering
\textbf{（a）理想负荷工况结果}\\[3pt]
\begin{tabular}{@{}lr@{}}
\toprule
指标 & 数值 \\
\midrule
$Q_{\mathrm{pc}}$ (kW) & 17.05 \\
$q_{\mathrm{pc}}$ (W/m$^{2}$) & 12.47 \\
$E_{\mathrm{cool}}$ (kWh) & 13486 \\
$E_{\mathrm{heat}}$ (kWh) & 169458 \\
$E_{\mathrm{total}}$ (kWh) & 182944 \\
$E/A_{\mathrm{f}}$ (kWh/m$^{2}$) & 200.74 \\
\bottomrule
\end{tabular}
\end{minipage}%
\hfill%
\begin{minipage}[t]{0.54\textwidth}
\centering
\textbf{（b）两类系统比较}\\[3pt]
\begin{tabular}{@{}lcc@{}}
\toprule
指标 & VRF+DOAS & FCU+DOAS \\
\midrule
设计冷量 (kW) & 41.70 & 73.79 \\
设计热量 (kW) & 41.70 & 47.56 \\
年电耗 (kWh) & 54842 & 4381 \\
电耗强度 (kWh/m$^{2}$) & 60.20 & 4.81 \\
\bottomrule
\end{tabular}
\end{minipage}
\end{table}

天津理想负荷结果的显著特点是：峰值冷负荷为\SI{17.05}{\kilo\watt}，明显高于沈阳（\SI{10.00}{\kilo\watt}），但低于南方城市。年热负荷为\SI{169458}{\kilo\watt\hour}，仍占年总负荷的约92.6\%，表明供暖需求在天津仍占据主导地位。年冷负荷为\SI{13486}{\kilo\watt\hour}，约为年热负荷的8.0\%。

\subsection{空调系统方案比较}

天津VRF+DOAS的年电耗强度为\SI{60.20}{\kilo\watt\hour\per\square\meter}，略低于沈阳（\SI{60.25}{\kilo\watt\hour\per\square\meter}），反映出较温和的冬季使VRF热泵在供热工况下的运行时间有所减少。FCU+DOAS电耗强度为\SI{4.81}{\kilo\watt\hour\per\square\meter}，其年度天然气消耗为\SI{70174}{\kilo\watt\hour}（强度为\SI{77.03}{\kilo\watt\hour\per\square\meter}）。

天津FCU方案中，冷水机组设计冷量（\SI{73.79}{\kilo\watt}）与锅炉设计热量（\SI{47.56}{\kilo\watt}）之比约为1.55:1，与沈阳（冷热比约0.95:1，即54.71/57.38）形成对比，体现从北向南冷热需求量此消彼长的趋势。VRF方案的设计冷/热量均为\SI{41.70}{\kilo\watt}，位于五城市的中间水平。
